\section{De las moléculas a las células: la investigación en biología con IA a través de escalas}

\subsection{Un panorama más amplio que el del diseño de proteínas}

Amini enfatiza que el diseño de proteínas es solo un nivel dentro de este panorama más amplio de la IA en biología. La biología es un \textbf{sistema profundamente jerarquizado}:

\begin{table}[H]
\centering
\caption{La estructura jerárquica de la biología y los retos de IA correspondientes}
\begin{tabular}{@{}llll@{}}
\toprule
\textbf{Nivel} & \textbf{Objeto} & \textbf{Tipo de datos} & \textbf{Tarea de IA} \\
\midrule
Molecular & Proteínas, ADN & Secuencia, estructura & Diseño molecular, predicción de función \\
Celular & Estado de la célula individual & Transcriptoma, proteoma & Caracterización del estado celular, predicción de perturbaciones \\
Tisular & Cortes de tejido & Imágenes patológicas & Detección, clasificación y pronóstico de enfermedades \\
Individual & Datos del paciente & Multiómica, registros clínicos & Medicina de precisión, predicción de eficacia terapéutica \\
\bottomrule
\end{tabular}
\end{table}

La investigación del equipo MSR BioML abarca todos estos niveles, con el objetivo de desarrollar métodos de IA capaces de razonar en distintas escalas y, finalmente, integrarlos para impulsar decisiones con impacto clínico.

\subsection{La dirección de fusión entre secuencia y estructura}

En el campo del diseño de proteínas, el equipo está explorando activamente el \textbf{modelado conjunto} de la información de secuencia y de estructura. Amini señala que secuencia y estructura son complementarias:

\begin{itemize}
    \item La \textbf{estructura} proporciona un control geométrico preciso de alta resolución
    \item La \textbf{secuencia} proporciona una cobertura más amplia de la diversidad biológica
\end{itemize}

Entre las vías para combinar ambas están: mapear la secuencia y la estructura a un espacio de representación común, y entrenar conjuntamente modelos generativos de secuencia-estructura, entre otras.

\subsection{Multimodalidad texto-proteína}

Otra dirección de vanguardia es la \textbf{alineación multimodal} entre proteínas y lenguaje natural:

\begin{itemize}
    \item Dada una secuencia/estructura de proteína $\rightarrow$ generar un texto de descripción funcional
    \item Dado un texto de descripción funcional $\rightarrow$ diseñar una proteína que satisfaga esa descripción
\end{itemize}

Esto avanza hacia la visión que Amini describió al inicio de la conferencia: usar un prompt en lenguaje natural (como «diseña una proteína capaz de dirigirse a células de cáncer de mama y unirse a ellas») para impulsar el diseño de proteínas.

\subsection{IA a escala celular: entender los estados de enfermedad}

A nivel celular, el equipo está desarrollando métodos de IA capaces de definir y caracterizar los \textbf{estados celulares} (cellular states). El objetivo es entender cómo la enfermedad (en particular el cáncer) modifica el estado celular, y usar esa información para orientar la selección de estrategias terapéuticas.

\subsection{Patología digital: análisis de imágenes de tejido impulsado por IA}

A nivel tisular, el equipo de MSR ha logrado avances decisivos en el campo de la \textbf{patología digital} (digital pathology): la aplicación de métodos de visión por computadora al análisis de imágenes histopatológicas de tejido, para lograr la detección, clasificación y pronóstico automáticos del cáncer.

\begin{knowledgebox}{La colaboración estratégica entre MSR y el Broad Institute}
El equipo MSR BioML ha establecido una colaboración de investigación estratégica y de largo plazo con el Broad Institute del MIT y Harvard. El Broad Institute es una institución líder a nivel mundial en investigación de medicina de precisión y genómica. El objetivo de esta colaboración es articular estrechamente las capacidades de IA con recursos experimentales y clínicos de primer nivel, para lograr un ciclo completo que va de los datos moleculares del paciente a la predicción de IA, de ahí a la validación experimental y, finalmente, a la aplicación clínica. Este modelo de colaboración refleja el rasgo central del campo de la IA para la biología: la imprescindible complementariedad entre lo computacional y lo experimental.
\end{knowledgebox}

\subsection{Resumen de la sección}

La investigación en IA para la biología va mucho más allá del diseño de proteínas: es un proyecto sistemático que atraviesa múltiples niveles, desde el molecular hasta el celular, el tisular y el individual. El trabajo del equipo MSR BioML cubre desde el modelado conjunto de secuencia y estructura de proteínas hasta la patología digital, con la meta final de integrar métodos de IA multiescala para impulsar decisiones de medicina de precisión con impacto real en la práctica clínica.
